\capitulo{1}{Introducción}
Durante las últimas décadas se ha observado un aumento considerable en el número de afectados por la diabetes. Las previsiones para 2035 coinciden en que la cifra de diabéticos se duplicará hasta alcanzar los 5,1 millones de afectados \cite{idf_atlas_diabetes}. Para los pacientes, esto supone un cambio drástico en su estilo de vida y la necesidad de un seguimiento constante. Por ello, resulta vital mantener un control exhaustivo de la enfermedad y fomentar el autoconocimiento.

Debido a las exigencias y cuidados requeridos, un control inadecuado puede derivar en la aparición de patologías asociadas. Estas complicaciones pueden afectar gravemente a órganos vitales como el corazón o los riñones, así como provocar problemas circulatorios en las extremidades o causar enfermedades oculares severas como las retinopatías.

Resulta fundamental que las personas diagnosticadas sepan gestionar su condición y dispongan de las herramientas tecnológicas adecuadas. En la actualidad existen múltiples dispositivos médicos, tales como glucómetros, sensores de medición continua de glucosa, bombas de insulina o bolígrafos inteligentes, entre otros. Sin embargo, en el panorama actual no existe un ecosistema unificado donde el usuario pueda centralizar la información y aprovechar al máximo las ventajas conjuntas que ofrecen estas herramientas. 

Adicionalmente, uno de los puntos más vulnerables en el tratamiento es la educación diabetológica. Esta responsabilidad suele estar delegada en un reducido número de especialistas dentro del entorno hospitalario. Esto provoca que el seguimiento continuado, especialmente en entornos rurales, resulte complejo y menos accesible, forzando a que el control diario y la toma de decisiones deban ser asumidos casi en su totalidad por el propio paciente.

En este contexto nace SugAIr, un proyecto cuyo objetivo es el desarrollo de un asistente conversacional inteligente diseñado para ofrecer soporte y educación diabetológica de forma accesible y privada. Mediante una aplicación web, el sistema integra Inteligencia Artificial Generativa de ejecución local. Esta decisión arquitectónica permite resolver dudas en tiempo real garantizando que la información médica sensible del usuario no abandone su propio entorno de ejecución.

Todo el código fuente del proyecto, así como el histórico de desarrollo y las versiones de los modelos empleados, están disponibles públicamente en el repositorio de GitHub del proyecto: 
\url{https://github.com/igorArroyo/TFG-de-Igor}

Del mismo modo, el seguimiento del desarrollo, la definición de iteraciones y la gestión ágil del trabajo se han administrado mediante un tablero Kanban, disponible en la plataforma GitLab:
\url{https://gitlab.com/HP-SCDS/Observatorio/2025-2026/sugair/ubu-sugair/-/boards}
(No soy el administrador de este repositorio, el que quiera acceso que me lo pida y hablo con Luis para que os de acceso)


